\documentclass[10pt,letterpaper]{article}

\usepackage[margin=0.48in]{geometry}
\usepackage[T1]{fontenc}
\usepackage{lmodern}
\usepackage{xcolor}
\usepackage{enumitem}
\usepackage{tabularx}
\usepackage[hidelinks]{hyperref}
\usepackage{titlesec}
\usepackage{microtype}

\definecolor{accent}{HTML}{8A3C52}
\setlength{\parindent}{0pt}
\setlength{\parskip}{0pt}
\setlist[itemize]{leftmargin=1.15em,itemsep=0.7pt,topsep=1pt,parsep=0pt}
\titleformat{\section}{\large\bfseries\color{accent}}{}{0pt}{}[\vspace{-2pt}\titlerule]
\titlespacing*{\section}{0pt}{4pt}{2pt}
\pagestyle{empty}

\newcommand{\role}[4]{%
  \textbf{#1}\hfill #2\\
  \textit{#3}\hfill \textit{#4}\vspace{-3pt}
}

\begin{document}
\fontsize{9}{10.4}\selectfont

\begin{center}
  {\Large\bfseries Piter Z. Garcia Bautista}\\[1pt]
  {\large Inclusive Computer Science Education, Accessibility, and Applied AI}\\[2pt]
  Rochester, NY\\
  \href{mailto:pgarcia8@u.rochester.edu}{pgarcia8@u.rochester.edu}
  \enspace|\enspace +1 (646) 288-3300
  \enspace|\enspace \href{https://linkedin.com/in/garciapiterz}{linkedin.com/in/garciapiterz}\\
  \href{https://github.com/pzg8794}{github.com/pzg8794}
  \enspace|\enspace \href{https://garciapiterz.com}{garciapiterz.com}
\end{center}

\section*{Profile}
Graduate educator and applied AI researcher connecting inclusive computer
science teaching, disability access, responsible technology, and multilingual
communication. Current preparation includes NSF-funded K--12 computer science
teacher education, graduate work in data science and AI, an advanced
certificate in disability and inclusive practice, and professional experience
building clear, reproducible data systems for varied audiences.

\section*{Education}
\textbf{M.S., Teaching Computer Science K--12}, University of Rochester
\hfill Expected December 2026\\
GPA: 3.9+; NSF Robert Noyce Scholar; Advanced Certificate in Leadership in
Disability and Inclusive Practices

\textbf{M.S., Data Science}, Rochester Institute of Technology
\hfill Expected December 2026\\
GPA: 3.9+; applied AI, bioinformatics, neural networks, and reproducible data systems

\textbf{M.S., Computer Science}, Rochester Institute of Technology
\hfill 2015\\
Concentrations and project work in artificial intelligence, big-data
analytics, image processing, and computer graphics

\textbf{B.S., Computer Engineering Technology}, Farmingdale State College
\hfill 2012

\section*{Teaching, Research, and Professional Experience}
\role{Computer Science Teacher Candidate}{2025--Present}
{University of Rochester / Pine Brook Elementary School}{Rochester, NY}
\begin{itemize}
  \item Co-design standards-aligned learning in coding, robotics, digital
    fluency, and responsible AI for varied elementary learners.
  \item Use Universal Design for Learning, multimodal participation, and clear
    technical communication to expand access without lowering expectations.
\end{itemize}

\role{Graduate Assistant, Quantum and AI Research}{Fall 2025--Present}
{Rochester Institute of Technology, Software Engineering}{Rochester, NY}
\begin{itemize}
  \item Build reproducible Python evaluation workflows for adaptive
    decision-making under noisy, resource-constrained, and adversarial
    conditions.
  \item Translate complex methods and results into reports, presentations,
    teaching-ready explanations, and manuscript artifacts.
\end{itemize}

\role{Data Solutions Engineer}{Sep 2022--Aug 2024}
{VEDADATA Inc.}{New York, NY}
\begin{itemize}
  \item Designed Python data workflows, validation checks, and documentation;
    collaborated across technical and nontechnical teams to clarify needs and
    improve reliability and usability.
\end{itemize}

\role{AI Data Solutions Engineer}{Jan 2021--Sep 2022}
{VIOME Inc.}{New York, NY}
\begin{itemize}
  \item Developed healthcare-oriented machine-learning workflows for
    preprocessing, feature engineering, model experimentation, evaluation, and
    accessible reporting.
\end{itemize}

\section*{Selected Work}
\textbf{Inclusive CS and Disability Leadership.}
Develop teaching artifacts that use multiple means of engagement,
representation, and action while preserving rigorous computing goals.

\textbf{EQUITAS and the Puzzle Plan.}
Use EQUITAS, the equity-focused lens within the interdisciplinary Puzzle Plan,
to connect responsible data systems, diagnostic access, disability justice,
and inclusive education.

\textbf{Reproducible AI Evaluation.}
Build structured data-collection, validation, experiment-tracking, and
comparison workflows that make technical decisions inspectable and reusable.

\section*{Skills and Languages}
\begin{tabularx}{\textwidth}{@{}>{\bfseries}l X@{}}
Teaching and access & Inclusive CS curriculum, UDL, differentiated participation, assessment design, technical communication\\
Technical & Python, SQL, R, Java, C++, machine learning, data-quality validation, reproducible research, Git/GitHub\\
Languages & Spanish: native/fluent; English: fluent\\
\end{tabularx}

\section*{Recognition}
NSF Robert Noyce Scholar (2024--2026)
\enspace|\enspace IEEE Alumni Member (2009--Present)
\enspace|\enspace RIT/MESCYT merit-based graduate funding

\end{document}
